\chapter{Getting started}\label{ch:getting-started}

\section{Requirements and installation}

\app\ is a Python program. It needs:

\begin{itemize}
  \item Python 3.10 or later (it is tested with the Anaconda distribution, Python 3.12);
  \item \textbf{PySide6} (Qt~6) for the window;
  \item \textbf{numpy}, \textbf{scipy} and \textbf{networkx} for the engine.
\end{itemize}

Anaconda already ships numpy, scipy and networkx. PySide6 comes from pip:

\begin{terminal}
pip install PySide6 numpy scipy networkx
\end{terminal}

There is nothing to build. Get the repository (\cli{git clone} or a downloaded archive),
open a terminal in its root folder, the one that holds \file{run\_gui.py}, and start the
program from there. The console tools of \cref{ch:console} do not need PySide6.

\begin{note}
On Windows with Anaconda, Qt~6 can fail to load with \emph{``the specified procedure
could not be found''}, because Anaconda ships its own ICU library and that one shadows the
library Qt expects. \app\ works around it: the package \code{psa.gui} preloads the system
ICU when it is imported. \file{run\_gui.py} already imports it first. If you write a
script that uses PySide6 together with \app, import \code{psa.gui} before anything else
from PySide6.
\end{note}

\section{Starting the program}

\begin{terminal}
python run_gui.py                   # an empty project
python run_gui.py --example         # the validated Assignment S1 network (A = 80)
python run_gui.py my_network.psn    # open a project
\end{terminal}

\begin{itemize}
  \item \cli{--example} opens \file{examples/s1\_a80.psn} as an \emph{untitled copy}:
    \keys{\ctrl + S} asks where to save it and never overwrites the example.
    \menu{File > Load Assignment S1 example} does the same from inside the program.
  \item If the file named on the command line does not exist or is damaged, \app\ still
    starts, with an empty project, and tells you why it could not open the file.
  \item If the project links a variable script (\cref{sec:script}), \app\ asks before it
    runs the script. Opening a file never executes code by itself.
  \item The list of recent files and the window layout are stored per user. On Windows
    they live in the registry, under \file{HKCU\textbackslash Software\textbackslash
    PowerLab\textbackslash PowerLab}.
\end{itemize}

\section{A first session in five minutes}\label{sec:first-session}

This walk-through uses the load flow example that comes with \app, so you can see the
whole cycle before drawing anything yourself.

\begin{steps}
  \item Start \app\ and open \file{examples/lf\_220kv.psn} with \menu{File > Open...}
    (\keys{\ctrl + O}). The canvas shows eight 220~kV buses. The slack generator G1 is at
    bus~1, an SVC at bus~3, generator G7 at bus~7, constant-power loads at buses 4, 6 and
    8, and a capacitor at bus~5.
  \item Look at the \gui{Study cases} panel on the right. The project holds three cases:
    \gui{LFa-LFc) Load flow}, \gui{LFd)}, which raises one load by 10\,\%, and the
    \gui{LF variant)} with the shunt rated as a susceptance.
  \item Press \keys{F5} (\menu{Calculate > Run all cases}). Each case gets a tab in the
    \gui{Results} panel at the bottom, and the diagram is annotated with the first one.
    Every bus shows its voltage, angle and type, every line its power flow and losses,
    and every generator its $P$ and $Q$.
  \item Read the headline of the first tab: \emph{Converged in 4 iterations. P\_loss =
    22.1658 MW}. Scroll down for the bus table, the generation table, the branch flows,
    the power balance, the iteration log and the Jacobian.
  \item Click the \gui{LFd)} tab. The annotations switch to that case: the losses rise to
    24.3993~MW.
  \item In the \gui{Variables} panel, double-click the \gui{Expression} cell of \code{A},
    change \code{80.0} to \code{0} and press \keys{Enter}. The loads that depend on \code{A} change. The result
    tabs are marked with an asterisk, and a banner over the canvas says that the results
    are out of date. Press \keys{F5} again: the losses are now 26.4240~MW, the answer
    printed in the assignment for $A=0$.
  \item Try \menu{View > Colour by > Voltage result} to colour each bus by its voltage
    band, and \menu{File > Export diagram...} to save the annotated drawing as a PNG, SVG
    or PDF file.
\end{steps}

\begin{tip}
\menu{Edit > Undo} (\keys{\ctrl + Z}) takes back the change of \code{A} too. Undo covers
every change to the network, the layout, the variables, the bases and the study cases.
\end{tip}

\section{The example projects}

The folder \file{examples/} holds one project per assignment. All of them are parametric
in the variable \code{A}, the student's A-number, and are set to $A=80$. Set it to $0$ to
reproduce the answers printed in the assignment statements.

\begin{reftable}{@{}>{\raggedright\arraybackslash}p{3.8cm}X@{}}
\toprule
\sffamily\bfseries File & \sffamily\bfseries Contents \\ \midrule
\file{s1\_a80.psn} & Assignment S1: a factory fed from a 40~kV grid through two
  transformers and two cable sections. Cases for the per-unit model, S1a, S1b and S1c. \\
\file{s2\_grid.psn} & Assignment S2 and S3a--S3c: the meshed 40~kV grid fed from 70~kV,
  with its zero-sequence data. \\
\file{grid\_industry.psn} & The S2 grid with the S1 industry connected at bus~7: cases
  S2a--S2d and S3a--S3e. \\
\file{lf\_220kv.psn} & Assignment LF: the eight-bus 220~kV network for the load flow,
  with cases LFa--LFd and the constant-B variant. \\
\bottomrule
\end{reftable}

\Cref{ch:assignments} walks through every task of these projects, with the case that
answers it and the expected numbers.
